% ===== Topic 2, part d: problems 41-50 (EN) =====
\providecommand{\tbGrid}[4]{\begin{tikzpicture}[baseline=0.2cm,x=0.5cm,y=0.5cm]
  \draw[Ink,line width=0.7pt] (0,0) grid[step=1] (2,2);
  \node[font=\sffamily\bfseries,text=tmBlue!80!black,inner sep=0pt,#4] at (#1,#2) {#3};\end{tikzpicture}}
% helper inside a tikzpicture: \tbG{x0}{y0}{cx}{cy}{letter}{transform} draws a 2x2 grid (cell 0.5) at (x0,y0)
\providecommand{\tbG}[6]{\draw[Ink,line width=0.7pt] (#1,#2) grid[step=0.5] ++(1,1);
  \node[font=\sffamily\bfseries,text=tmBlue!80!black,inner sep=0pt,#6] at ({#1+#3},{#2+#4}) {#5};}

% ---------------------------------------------------------------- 41
\begin{probbox}{41}{\stars{1}}
\textit{(Extension of Problem 16.)} We only consider rectangles whose two side lengths are \emph{positive
integers} (in cm); a square also counts as a rectangle.
\begin{enumerate}[label=\alph*),leftmargin=1.6em,itemsep=1pt]
\item Find all rectangles whose area (in cm$^2$) is numerically \emph{equal} to their perimeter (in cm).
\item Find all rectangles whose area is numerically \emph{twice} their perimeter.
\end{enumerate}
\end{probbox}
\probsol{
\textbf{Understanding the problem.} Call the sides $a$ and $b$ (positive integers). Without loss of generality,
assume $a\le b$ ($a$ is the shorter side). The area is $ab$ and the perimeter is $2(a+b)$. This is a ``whole-number
solutions'' problem: one equation with two unknowns. Trick: use $a\le b$ to \emph{trap} $a$ in a small range, then
try each value.

\textbf{a)} Condition: $ab=2a+2b$.

\step{1}{An upper bound for $a$}
Since $a\le b$, we have $2a+2b\le 2b+2b=4b$. Hence $ab\le 4b$. Dividing both sides by $b>0$: $a\le 4$.

\step{2}{A lower bound for $a$}
Since $2a>0$, we have $ab=2a+2b>2b$. Dividing both sides by $b>0$: $a>2$. So $a=3$ or $a=4$.

\step{3}{Try each value}
\begin{itemize}[leftmargin=1.4em,itemsep=1pt]
\item $a=3$: $3b=6+2b\Rightarrow b=6$. The $3\times 6$ rectangle.
\item $a=4$: $4b=8+2b\Rightarrow 2b=8\Rightarrow b=4$. The $4\times 4$ square.
\end{itemize}
\begin{figC}\begin{scope}[x=0.45cm,y=0.45cm]
\fill[tmBlue!15] (0,0) rectangle (6,3);\draw[Line] (0,0) grid[step=1] (6,3);\draw[tmBlue!80!black,line width=1pt] (0,0) rectangle (6,3);
\node[below,font=\footnotesize] at (3,0){$A=18$, $P=18$};
\begin{scope}[xshift=4cm]
\fill[tmGreen!15] (0,0) rectangle (4,4);\draw[Line] (0,0) grid[step=1] (4,4);\draw[tmGreen!70!black,line width=1pt] (0,0) rectangle (4,4);
\node[below,font=\footnotesize] at (2,0){$A=16$, $P=16$};
\end{scope}
\end{scope}\end{figC}

\textbf{b)} Condition: $ab=2\cdot 2(a+b)=4a+4b$.

\step{4}{Trap $a$ as before}
$ab=4a+4b\le 8b\Rightarrow a\le 8$; and $ab=4a+4b>4b\Rightarrow a>4$. So $a\in\{5,6,7,8\}$.

\step{5}{Try each value}
From $ab=4a+4b$ we get $ab-4b=4a$, i.e.\ $b(a-4)=4a$, so $b=\dfrac{4a}{a-4}$.
\begin{center}\begin{tabular}{|c|c|c|c|c|}\hline
$a$ & $5$ & $6$ & $7$ & $8$\\\hline
$b=\dfrac{4a}{a-4}$ & $\dfrac{20}{1}=20$ & $\dfrac{24}{2}=12$ & $\dfrac{28}{3}$ (not an integer) & $\dfrac{32}{4}=8$\\\hline
\end{tabular}\end{center}
There are three rectangles: $5\times 20$, $6\times 12$, $8\times 8$.

\begin{proofbox}Why have we missed none? Every rectangle that works has its shorter side $a$ in the range we
trapped, and for each $a$ the side $b$ is determined uniquely by the equation. We tried all those values of
$a$.\end{proofbox}
\begin{tipbox}Another way for part a): $ab-2a-2b+4=4$, i.e.\ $(a-2)(b-2)=4=1\cdot 4=2\cdot 2$. Hence
$a-2=1,\ b-2=4$ or $a-2=b-2=2$. Again we get $3\times 6$ and $4\times 4$.\end{tipbox}
\verify{$3\times 6$: $A=18$, $P=2\cdot 9=18$. $4\times4$: $16=16$. $5\times 20$: $A=100$, $P=50$;
$6\times 12$: $A=72$, $P=36$; $8\times 8$: $A=64$, $P=32$ -- all twice.}
\ans{a) $3\times 6$ and $4\times 4$. \quad b) $5\times 20$, $6\times 12$, $8\times 8$.}
}

% ---------------------------------------------------------------- 42
\begin{probbox}{42}{\stars{1}}
\textit{(Extension of Problem 17.)} The letter P is written in the top-right cell of a $2\times 2$ grid:
\tbGrid{1.5}{1.5}{P}{}\,. We may use three transformations, as many times as we like and in any order:
\begin{itemize}[leftmargin=1.4em,itemsep=1pt]
\item $Q$: rotate the grid $90^\circ$ clockwise about its centre;
\item $D$: reflect the grid in the \emph{vertical} line through the centre;
\item $N$: reflect the grid in the \emph{horizontal} line through the centre.
\end{itemize}
\begin{enumerate}[label=\alph*),leftmargin=1.6em,itemsep=1pt]
\item Draw the result of doing $D$ and then $N$. Can that result be obtained using only $Q$?
\item Show that doing $D$ and then $Q$ has the same effect as a reflection in one of the \emph{diagonals} of the
grid. Which diagonal?
\item How many different pictures can be obtained altogether (including the starting picture)? Explain why there
cannot be more.
\end{enumerate}
\end{probbox}
\probsol{
\textbf{Understanding the problem.} As in Problem 17, think of the grid as a clear square \emph{card} with the
letter P drawn on it. Each transformation lifts the card and puts it back exactly in its place. To keep track we
record two things: which cell P is in, and how it has been turned or flipped.

\textbf{a)}
\step{1}{Do $D$}
Reflecting in the vertical line: the top-right cell goes to the top-left cell, and P is ``mirrored'' (its bowl
now faces left).
\step{2}{Then do $N$}
Reflecting in the horizontal line: the top-left cell goes down to the bottom-left cell, and the letter is flipped
top-to-bottom. Result: P is in the bottom-left cell, \emph{upside down} (looking like a ``d''), as in the third
picture below.
\begin{figC}
\tbG{0}{0}{0.75}{0.75}{P}{}\node at (1.75,0.5){$\xrightarrow{\;D\;}$};
\tbG{2.5}{0}{0.25}{0.75}{P}{xscale=-1}\node at (4.25,0.5){$\xrightarrow{\;N\;}$};
\tbG{5}{0}{0.25}{0.25}{P}{rotate=180}
\node[anchor=west] at (6.2,0.5){\small$=$ rotation by $180^\circ$};
\end{figC}
\step{3}{Compare with the $Q$ moves}
Doing $Q$ twice (a $180^\circ$ rotation): the top-right cell goes to the opposite cell through the centre, the
bottom-left cell, and the letter is upside down. This is the same result. So \textbf{$D$ then $N$ has the same
effect as $Q$ then $Q$.}
\begin{recallbox}Two reflections in two perpendicular lines (both through the centre) have the same effect as a
$180^\circ$ rotation about the centre.\end{recallbox}

\textbf{b)}
\step{4}{Do $D$ and then $Q$}
$D$: P goes to the top-left cell, mirrored. $Q$ ($90^\circ$ clockwise): the top-left cell goes to the top-right
cell. Result: P is \emph{back in the top-right cell}, but its stem is horizontal and its bowl is on top:
\begin{figC}
\tbG{0}{0}{0.75}{0.75}{P}{}\node at (1.75,0.5){$\xrightarrow{\;D\;}$};
\tbG{2.5}{0}{0.25}{0.75}{P}{xscale=-1}\node at (4.25,0.5){$\xrightarrow{\;Q\;}$};
\tbG{5}{0}{0.75}{0.75}{P}{cm={0,1,1,0,(0,0)}}\draw[tmOrange,dashed,line width=0.9pt] (4.85,-0.15)--(6.15,1.15);
\draw[tmOrange,dashed,line width=0.9pt] (4.05,-0.5)--(5.15,0.6);
\end{figC}
\step{5}{Recognise the diagonal reflection}
The top-right cell has not moved, which means it lies on the mirror line. The diagonal through the top-right cell
(and the bottom-left cell) is the diagonal joining the \textbf{bottom-left corner to the top-right corner}
(dashed). Reflecting in this line turns ``up'' into ``right'' and vice versa, so the (vertical) stem of P becomes
horizontal and the bowl of P (on the right) moves to the top -- exactly as in the picture. A further check: a
letter in the top-left cell moves to the bottom-right cell in both cases.

\textbf{c)}
\step{6}{There cannot be more than $8$ pictures}
\begin{proofbox}Mark the top-right corner of the card with a star. Each time the card is put back, the star can
only be in $1$ of the $4$ corners of the frame. In addition, the card can only be \emph{face up} or \emph{face
down} (each reflection $D$ or $N$ switches face up and face down; the rotation $Q$ does not). Once we know where
the star is and whether the card is face up or face down, the position of the card -- and hence the whole
picture -- is completely determined. So there are at most $4\times 2=8$ pictures.\end{proofbox}
\step{7}{All $8$ pictures occur}
Four ``face up'' pictures: rotations by $0^\circ$, $90^\circ$, $180^\circ$, $270^\circ$ (using $Q$ from $0$ to $3$
times). Four ``face down'' pictures: each of the above followed by $D$. The letter P has no line of symmetry, so
these $8$ pictures are all different:
\begin{figC}
\tbG{0}{0}{0.75}{0.75}{P}{}\tbG{1.5}{0}{0.75}{0.25}{P}{rotate=-90}
\tbG{3}{0}{0.25}{0.25}{P}{rotate=180}\tbG{4.5}{0}{0.25}{0.75}{P}{rotate=90}
\tbG{0}{-1.5}{0.25}{0.75}{P}{xscale=-1}\tbG{1.5}{-1.5}{0.25}{0.25}{P}{cm={0,-1,-1,0,(0,0)}}
\tbG{3}{-1.5}{0.75}{0.25}{P}{yscale=-1}\tbG{4.5}{-1.5}{0.75}{0.75}{P}{cm={0,1,1,0,(0,0)}}
\node[anchor=east,font=\footnotesize] at (-0.2,0.5){face up};\node[anchor=east,font=\footnotesize] at (-0.2,-1){face down};
\end{figC}
\begin{warnbox}Do not count ``$4$ cells $\times$ $8$ orientations $=32$''. The cell and the orientation of the
letter are \emph{tied together}, because both move with the same card.\end{warnbox}
\ans{a) P in the bottom-left cell, upside down; the same as $Q$ then $Q$ (a $180^\circ$ rotation). \ b) The
diagonal joining the bottom-left corner to the top-right corner. \ c) $8$ pictures.}
}

% ---------------------------------------------------------------- 43
\begin{probbox}{43}{\stars{1}}
\textit{(Extension of Problem 18.)} Rectangle $ABCD$ has $AB=12$ cm and $BC=8$ cm. Point $P$ lies inside the
rectangle. (We write $S_{XYZ}$ for the area of triangle $XYZ$.)
\begin{figC}\begin{scope}[x=0.36cm,y=0.36cm]
\fill[tmBlue!25] (0,0)--(12,0)--(7,3)--cycle;\fill[tmOrange!25] (12,8)--(0,8)--(7,3)--cycle;
\draw[Ink,line width=1pt] (0,0) rectangle (12,8);
\draw[Ink] (0,0)--(7,3)--(12,0) (0,8)--(7,3)--(12,8);
\fill[Ink] (7,3) circle (2.5pt);\node[right=2pt] at (7,3){$P$};
\node[below left] at (0,0){$A$};\node[below right] at (12,0){$B$};\node[above right] at (12,8){$C$};\node[above left] at (0,8){$D$};
\end{scope}\end{figC}
\begin{enumerate}[label=\alph*),leftmargin=1.6em,itemsep=1pt]
\item Prove that wherever $P$ is, we always have $S_{PAB}+S_{PCD}=48$ cm$^2$.
\item Given $S_{PAB}=18$ cm$^2$ and $S_{PBC}=20$ cm$^2$, find $S_{PCD}$, $S_{PDA}$ and the distances from $P$ to
the four sides of the rectangle.
\item Is there a point $P$ inside the rectangle with $S_{PAB}=20$, $S_{PBC}=30$ and $S_{PCD}=30$ (cm$^2$)?
\end{enumerate}
\end{probbox}
\probsol{
\textbf{Understanding the problem.} Each of the four triangles has a side of the rectangle as its base and $P$ as
its apex. The height of each triangle is just \emph{the distance from $P$ to that side}. Call the distances from
$P$ to $AB$, $BC$, $CD$, $DA$ respectively $h_1,h_2,h_3,h_4$.
\begin{figC}\begin{scope}[x=0.36cm,y=0.36cm]
\draw[Ink,line width=1pt] (0,0) rectangle (12,8);
\draw[tmBlue,thick,<->] (7,0.1)--(7,2.9);\draw[tmOrange!80!black,thick,<->] (7,3.1)--(7,7.9);
\draw[tmGreen!70!black,thick,<->] (7.1,3)--(11.9,3);\draw[tmPurple,thick,<->] (0.1,3)--(6.9,3);
\fill[Ink] (7,3) circle (2.5pt);
\node[right,font=\footnotesize] at (7,1.5){$h_1$};\node[right,font=\footnotesize] at (7,5.5){$h_3$};
\node[above,font=\footnotesize] at (9.5,3){$h_2$};\node[above,font=\footnotesize] at (3.5,3){$h_4$};
\node[below left] at (0,0){$A$};\node[below right] at (12,0){$B$};\node[above right] at (12,8){$C$};\node[above left] at (0,8){$D$};
\end{scope}\end{figC}

\textbf{a)}
\step{1}{The sum of two heights}
Draw through $P$ the line perpendicular to $AB$; since $AB\parallel CD$ it is also perpendicular to $CD$. The part
between $AB$ and $CD$ has length $BC=8$, and $P$ divides it into two parts $h_1$ and $h_3$. So $h_1+h_3=8$.
Similarly $h_2+h_4=AB=12$.
\step{2}{Add the two areas}
\begin{recallbox}Area of a triangle $=\tfrac12\times$ base $\times$ height.\end{recallbox}
\[ S_{PAB}+S_{PCD}=\frac{12\cdot h_1}{2}+\frac{12\cdot h_3}{2}=6(h_1+h_3)=6\cdot 8=48 \text{ (cm}^2). \]
The result does not depend on the position of $P$. (It is exactly half the area of the rectangle,
$12\cdot 8=96$.) Similarly, $S_{PBC}+S_{PDA}=\tfrac{8}{2}(h_2+h_4)=4\cdot 12=48$ cm$^2$.

\textbf{b)}
\step{3}{Use part a}
$S_{PCD}=48-18=30$ cm$^2$, \quad $S_{PDA}=48-20=28$ cm$^2$.
\step{4}{Compute the four distances}
From $A=\tfrac12\cdot$ base $\cdot h$ we get $h=2A\div$ base.
\[ h_1=\frac{2\cdot 18}{12}=3,\quad h_2=\frac{2\cdot 20}{8}=5,\quad h_3=\frac{2\cdot 30}{12}=5,\quad h_4=\frac{2\cdot 28}{8}=7 \text{ (cm)}. \]

\textbf{c)}
\step{5}{Check the compulsory condition}
By part a), we always have $S_{PAB}+S_{PCD}=48$. But $20+30=50\ne 48$. So \textbf{there is no} such point $P$.

\begin{tipbox}Problem 18 is a special case: $P$ is the centre of a square, so the four triangles are equal. Here $P$
is arbitrary, yet ``the sum of two opposite triangles'' is still always half of the figure.\end{tipbox}
\verify{$h_1+h_3=3+5=8=BC$ and $h_2+h_4=5+7=12=AB$. With $A(0,0)$ we get $P(7,3)$ -- exactly the point in the
picture.}
\ans{b) $S_{PCD}=30$ cm$^2$, $S_{PDA}=28$ cm$^2$; the distances to $AB,BC,CD,DA$ are $3,5,5,7$ cm. \ c) No.}
}

% ---------------------------------------------------------------- 44
\begin{probbox}{44}{\stars{1}}
\textit{(Extension of Problem 22.)} Parallelogram $ABCD$ has area $48$ cm$^2$.
\begin{enumerate}[label=\alph*),leftmargin=1.6em,itemsep=1pt]
\item Let $M$ and $N$ be the midpoints of $BC$ and $CD$ respectively. Find the area of triangle $AMN$.
\item Keep $N$ as the midpoint of $CD$, but take $M$ on side $BC$ with $BM=3\,MC$. Find the area of triangle $AMN$.
\end{enumerate}
\begin{figC}\begin{scope}[x=0.55cm,y=0.55cm]
\fill[tmBlue!25] (0,0)--(7,2)--(5,4)--cycle;
\draw[Ink,line width=1pt] (0,0)--(6,0)--(8,4)--(2,4)--cycle;
\draw[tmBlue!80!black,line width=1pt] (0,0)--(7,2)--(5,4)--cycle;
\node[below left] at (0,0){$A$};\node[below right] at (6,0){$B$};\node[above right] at (8,4){$C$};\node[above left] at (2,4){$D$};
\node[right] at (7,2){$M$};\node[above] at (5,4){$N$};
\end{scope}\end{figC}
\end{probbox}
\probsol{
\textbf{Understanding the problem.} Triangle $AMN$ sits ``diagonally'' inside the parallelogram. We take the whole
parallelogram and remove the three corner triangles $ABM$, $MCN$ and $NDA$ (as usual, $S_{XYZ}$ denotes the area
of triangle $XYZ$):
\[ S_{AMN}=S_{ABCD}-S_{ABM}-S_{MCN}-S_{NDA}. \]
\begin{recallbox}(1) Each diagonal divides a parallelogram into two triangles of equal area ($=24$ cm$^2$).
(2) Two triangles with a common apex whose bases lie on the same line have areas in the same ratio as their bases.
\end{recallbox}

\textbf{a)}
\step{1}{Triangle $ABM$}
Compared with triangle $ABC$ ($=24$): common apex $A$, base $BM=\tfrac12 BC$. So $S_{ABM}=\tfrac12\cdot 24=12$.
\step{2}{Triangle $NDA$}
Compared with triangle $ACD$ ($=24$): common apex $A$, base $DN=\tfrac12 DC$. So $S_{NDA}=12$.
\step{3}{Triangle $MCN$ (in two steps, as in Problem 22)}
Triangle $BCD$ has area $24$. Triangle $MCD$ shares the apex $D$ with it, with base $CM=\tfrac12 CB$:
$S_{MCD}=12$. Triangle $MCN$ shares the apex $M$ with triangle $MCD$, with base $CN=\tfrac12 CD$:
$S_{MCN}=\tfrac12\cdot 12=6$.
\step{4}{Result}
$S_{AMN}=48-12-6-12=18$ cm$^2$.

\textbf{b)}
\step{5}{Redo it with $MC=\tfrac14 BC$}
Since $BM=3\,MC$, we have $BC=4\,MC$, i.e.\ $MC=\tfrac14 BC$ and $BM=\tfrac34 BC$.
\begin{itemize}[leftmargin=1.4em,itemsep=1pt]
\item $S_{ABM}=\tfrac34\cdot 24=18$;
\item $S_{NDA}=12$ (unchanged);
\item $S_{MCD}=\tfrac14\cdot 24=6$, then $S_{MCN}=\tfrac12\cdot 6=3$.
\end{itemize}
$S_{AMN}=48-18-3-12=15$ cm$^2$.

\begin{tipbox}The ``corner'' rule: if on two sides of a triangle, measured from the same vertex, we take the parts
$\tfrac{p}{q}$ and $\tfrac{r}{s}$ of those sides, then the small triangle at that corner has area
$\tfrac{p}{q}\cdot\tfrac{r}{s}$ of the big triangle. In part b): $\tfrac14\cdot\tfrac12=\tfrac18$ of $24$ is
$3$.\end{tipbox}
\verify{Put $A(0,0)$, $B(8,0)$, $D(2,6)$, $C(10,6)$ (area $8\cdot 6=48$). Part a): $M(9,3)$, $N(6,6)$; enclose
triangle $AMN$ in a $9\times 6=54$ rectangle and subtract $\tfrac{9\cdot3}{2}=13.5$, $\tfrac{3\cdot 3}{2}=4.5$,
$\tfrac{6\cdot6}{2}=18$: $18$ remains. Correct.}
\ans{a) $18$ cm$^2$. \quad b) $15$ cm$^2$.}
}

% ---------------------------------------------------------------- 45
\begin{probbox}{45}{\stars{1}}
\textit{(Extension of Problem 21.)} Consider triangles whose three side lengths are positive integers and whose
perimeter is $15$ (two triangles with the same three side lengths count as one).
\begin{enumerate}[label=\alph*),leftmargin=1.6em,itemsep=1pt]
\item Let $c$ be the longest side. Prove that $5\le c\le 7$.
\item How many such triangles are there? List them.
\item How many of them are isosceles? How many have \emph{exactly} two equal sides?
\end{enumerate}
\end{probbox}
\probsol{
\textbf{Understanding the problem.} Call the sides $a\le b\le c$ (in increasing order, so that nothing is counted
twice). We have $a+b+c=15$, and by the triangle inequality we only need to check $a+b>c$.
\begin{recallbox}Three segments $a\le b\le c$ are the sides of a triangle exactly when $a+b>c$ (the two smaller
sides add up to more than the largest one).\end{recallbox}

\textbf{a)}
\step{1}{Upper bound}
$a+b=15-c$. The condition $a+b>c$ becomes $15-c>c$, i.e.\ $2c<15$, $c<7.5$. Since $c$ is an integer, $c\le 7$.
\step{2}{Lower bound}
$c$ is the largest side, so $a\le c$ and $b\le c$, hence $15=a+b+c\le 3c$, i.e.\ $c\ge 5$.

\textbf{b)}
\step{3}{List them by $c$}
For each $c$ we have $a+b=15-c$ and $a\le b\le c$.
\begin{center}\begin{tabular}{|c|c|l|c|}\hline
$c$ & $a+b$ & Pairs $(a,b)$ with $a\le b\le c$ & Number of triangles\\\hline
$7$ & $8$ & $(1,7),\ (2,6),\ (3,5),\ (4,4)$ & $4$\\
$6$ & $9$ & $(3,6),\ (4,5)$ \quad (not $(2,7)$ since $7>6$) & $2$\\
$5$ & $10$ & $(5,5)$ & $1$\\\hline
\end{tabular}\end{center}
Altogether $4+2+1=7$ triangles:
\[ (1,7,7),\ (2,6,7),\ (3,5,7),\ (4,4,7),\ (3,6,6),\ (4,5,6),\ (5,5,5). \]
\begin{warnbox}For $c=7$, the pair $(1,7)$ gives the triangle $1,7,7$ -- very ``thin'' but still valid because
$1+7=8>7$. On the other hand, no triangle has a side $8$, because the other two sides would add up to
$7<8$.\end{warnbox}

\textbf{c)}
\step{4}{Pick out the isosceles triangles}
An isosceles triangle has at least two equal sides: $(1,7,7)$, $(4,4,7)$, $(3,6,6)$, $(5,5,5)$ -- that is $4$
triangles. Among them $(5,5,5)$ is equilateral (three equal sides), so $3$ triangles have exactly two equal sides.
\verify{Each triple adds up to $15$ and satisfies $a+b>c$: $8>7$, $8>7$, $8>7$, $8>7$, $9>6$, $9>6$, $10>5$.
Correct.}
\ans{a) $5\le c\le 7$. \ b) $7$ triangles. \ c) $4$ isosceles triangles, $3$ of which have exactly two equal sides.}
}

% ---------------------------------------------------------------- 46
\begin{probbox}{46}{\stars{2}}
\textit{(Extension of Problem 23.)} On a square grid (each square has side $1$), triangle $ABC$ has vertices
$A(0,0)$, $B(12,0)$, $C(3,9)$. Exactly $43$ points where gridlines cross lie \emph{inside} the triangle (not on its
sides). For each such point $P$ we get three triangles $PAB$, $PBC$, $PCA$.
\begin{figC}\begin{scope}[x=0.36cm,y=0.36cm]
\draw[Line,line width=0.5pt] (0,0) grid[step=1] (12,9);
\draw[Ink,line width=1.3pt] (0,0)--(12,0)--(3,9)--cycle;
\node[below left] at (0,0){$A$};\node[below right] at (12,0){$B$};\node[above] at (3,9){$C$};
\end{scope}\end{figC}
\begin{enumerate}[label=\alph*),leftmargin=1.6em,itemsep=1pt]
\item Find the area of triangle $ABC$.
\item Find the point $P$ for which the three triangles $PAB$, $PBC$, $PCA$ have equal areas.
\item Of all the $3\cdot 43=129$ triangles, how many have an area of exactly $\tfrac13$ of the area of triangle
$ABC$? For how many positions of $P$ is at least one of the three triangles like that?
\end{enumerate}
\end{probbox}
\probsol{
\textbf{Understanding the problem.} This is the ``one-third version'' of Problem 23. The key idea is the same: for a
triangle with a fixed base, the area depends only on the height, and points with the same height lie on a line
parallel to the base.

\textbf{a)}
\step{1}{Area of triangle $ABC$}
The base $AB=12$ lies on the $x$-axis and the height is the $y$-coordinate of $C$, which is $9$. Writing $S_{XYZ}$ for
the area of triangle $XYZ$:
$S_{ABC}=\tfrac{12\cdot 9}{2}=54$. One third is $18$.

\textbf{c)} (done first; part b) follows immediately)
\step{2}{Triangle $PAB$ with area $18$}
The base is $AB=12$, so the height must be $2\cdot 18\div 12=3$, i.e.\ $P$ lies on the horizontal line $y=3$.
On this line: side $AC$ (going ``$1$ right, $3$ up'') passes through $(1,3)$; side $BC$ (going ``$1$ left, $1$
up'' from $B$) passes through $(9,3)$. The points inside: $x=2,3,\dots,8$, i.e.\ $7$ points.

\step{3}{Triangle $PBC$ with area $18$}
\begin{proofbox}Take $E$ on $AB$ with $EB=\tfrac13 AB=4$, i.e.\ $E(8,0)$. Triangles $EBC$ and $ABC$ share the apex
$C$, and the base $EB=\tfrac13 AB$, so $S_{EBC}=\tfrac13\cdot 54=18$. Hence $S_{PBC}=18$ exactly when $P$ is as far
from $BC$ as $E$ is, i.e.\ when $P$ lies on the line through $E$ parallel to $BC$ (on the inside of the
triangle).\end{proofbox}
$BC$ goes in the direction ``$1$ left, $1$ up''. Starting from $E(8,0)$ in that direction: $(7,1)$, $(6,2)$,
$(5,3)$, $(4,4)$, $(3,5)$, then $(2,6)$ -- this point is on $AC$ (on $AC$ the $y$-coordinate is $3$ times the
$x$-coordinate), so it does not count. That gives $5$ points.

\step{4}{Triangle $PCA$ with area $18$}
Similarly, take $F$ on $AB$ with $AF=4$, i.e.\ $F(4,0)$; $P$ must lie on the line through $F$ parallel to $AC$
(direction ``$1$ right, $3$ up''): $(5,3)$, then $(6,6)$ -- which is on $BC$ (since $6+6=12$), so it does not
count. That gives $1$ point.
\begin{figC}\begin{scope}[x=0.36cm,y=0.36cm]
\draw[Line,line width=0.5pt] (0,0) grid[step=1] (12,9);
\draw[Ink,line width=1.3pt] (0,0)--(12,0)--(3,9)--cycle;
\draw[tmBlue,line width=1pt] (1,3)--(9,3);
\draw[tmOrange!90!black,line width=1pt] (8,0)--(2,6);
\draw[tmGreen!70!black,line width=1pt] (4,0)--(6,6);
\foreach \p in {(2,3),(3,3),(4,3),(6,3),(7,3),(8,3)}{\fill[tmBlue] \p circle (3pt);}
\foreach \p in {(7,1),(6,2),(4,4),(3,5)}{\fill[tmOrange!90!black] \p circle (3pt);}
\fill[AnsRed] (5,3) circle (4pt);
\node[below left] at (0,0){$A$};\node[below right] at (12,0){$B$};\node[above] at (3,9){$C$};
\node[below] at (8,0){$E$};\node[below] at (4,0){$F$};
\end{scope}\end{figC}

\step{5}{Count}
Number of triangles: $7+5+1=13$. The point $(5,3)$ lies on all three lines (the red point), so it was counted three
times. Number of different positions of $P$: $13-2=11$.

\textbf{b)}
\step{6}{Three equal triangles}
The three triangles have total area $54$, so they are equal when each has area $18$, i.e.\ $P$ must lie on
\emph{all three} lines we found. There is only one such point: $P(5,3)$.

\begin{tipbox}The point $(5,3)$ is the \emph{centroid} of the triangle: $x=\tfrac{0+12+3}{3}=5$,
$y=\tfrac{0+0+9}{3}=3$. You will study the centroid in detail in the chapter on triangles.\end{tipbox}
\verify{For $P(5,3)$: $S_{PAB}=\tfrac{12\cdot 3}{2}=18$. $S_{PCA}$: enclosing rectangle $5\times 9=45$, minus
$\tfrac{3\cdot 9}{2}=13.5$ (top-left corner), $\tfrac{2\cdot 6}{2}=6$ (top-right corner), $\tfrac{5\cdot 3}{2}=7.5$
(bottom-right corner): $18$ remains. So $S_{PBC}=54-18-18=18$. Correct.}
\ans{a) $54$. \ b) $P(5,3)$. \ c) $13$ triangles, at $11$ positions of $P$.}
}

% ---------------------------------------------------------------- 47
\begin{probbox}{47}{\stars{1}}
\textit{(Extension of Problem 25.)} Again we arrange $4\times 2$ rectangles in a ``staircase'' as in Problem 25,
but this time the top-left vertex of each rectangle (from the second one on) is placed on the bottom side of the
rectangle just above it, at distance $s$ from that rectangle's bottom-left vertex, where $s$ is $1$, $2$ or $3$.
(Problem 25 is the case $s=2$.)
\begin{figC}\begin{scope}[x=0.38cm,y=0.38cm]
\foreach \i in {0,1,2,3}{\draw[Ink,line width=0.9pt] (\i,-2*\i) rectangle (\i+4,-2*\i+2);}
\draw[tmBlue,<->] (0,-0.5)--(1,-0.5);\node[below,font=\footnotesize] at (0.5,-0.5){$s$};
\node[above] at (2,2){$4$};\node[left] at (0,1){$2$};
\foreach \k in {0,1,2}{\fill[Ink] (5.5+0.5*\k,-6.6-0.5*\k) circle (2.3pt);}
\end{scope}\end{figC}
\begin{enumerate}[label=\alph*),leftmargin=1.6em,itemsep=1pt]
\item For $s=1$ and $10$ rectangles, find the perimeter of the whole figure.
\item Find a formula for the perimeter in terms of the number of rectangles $n$ and the distance $s$.
\item Can we choose $s\in\{1,2,3\}$ and a number of rectangles $n$ so that the perimeter is exactly $100$? Find all
the ways.
\end{enumerate}
\end{probbox}
\probsol{
\textbf{Understanding the problem.} In Problem 25 we saw that the perimeter of a ``staircase'' equals the perimeter
of the smallest rectangle enclosing it. That argument works for every $s$ (the top-right boundary only goes right
and down; so does the bottom-left boundary). So we only need the width and the height of the whole figure.
\begin{recallbox}A ``staircase'' shape (whose boundary on each side only goes right/down) has the same perimeter as
its enclosing rectangle: $2\times(\text{width}+\text{height})$.\end{recallbox}

\textbf{b)} (done first, since part a) follows from it)
\step{1}{Width and height}
Each new rectangle is shifted $s$ to the right. The first has width $4$ and each later one adds $s$:
width $=4+(n-1)s$. Each new rectangle is $2$ lower: height $=2n$.
\step{2}{The formula}
\[ P=2\big(4+(n-1)s+2n\big). \]

\textbf{a)}
\step{3}{Substitute $s=1$, $n=10$}
Width $4+9\cdot 1=13$, height $20$; $P=2(13+20)=66$.

\textbf{c)}
\step{4}{Consider each value of $s$}
\begin{itemize}[leftmargin=1.4em,itemsep=1pt]
\item $s=1$: $P=2(4+n-1+2n)=2(3n+3)=6n+6$. $6n+6=100\Rightarrow 6n=94$: not divisible, rejected.
\item $s=2$: $P=2(4+2n-2+2n)=2(4n+2)=8n+4$. $8n+4=100\Rightarrow n=12$. Accepted.
\item $s=3$: $P=2(4+3n-3+2n)=2(5n+1)=10n+2$. $10n+2=100\Rightarrow 10n=98$: rejected.
\end{itemize}
There is only one way: $s=2$ and $n=12$.
\begin{warnbox}Remember that the height $2n$ does not depend on $s$; only the width changes. Many students wrongly
add in the hidden pieces inside the figure.\end{warnbox}
\verify{Part a) by direct counting: vertical sides $10\cdot 4=40$; horizontal sides on top $4+9\cdot 1=13$; at the
bottom $9\cdot 1+4=13$. Total $66$. Correct. Part c): $s=2$, $n=12$: $8\cdot 12+4=100$. Correct.}
\ans{a) $66$. \ b) $P=2\big(4+(n-1)s+2n\big)$. \ c) Only $s=2$, $n=12$.}
}

% ---------------------------------------------------------------- 48
\begin{probbox}{48}{\stars{1}}
\textit{(Extension of Problem 27.)} The figure below consists of $15$ identical small parallelograms arranged in $3$
rows of $5$. The small parallelogram in the very middle is shaded.
\begin{figC}\begin{scope}[x=0.5cm,y=0.5cm]
\fill[tmOrange!45] (5.4,1)--(7.6,1)--(8.6,2)--(6.4,2)--cycle;
\foreach \k in {0,...,5}{\draw[Ink,line width=0.9pt] (2.2*\k,0)--(2.2*\k+3,3);}
\foreach \j in {0,...,3}{\draw[Ink,line width=0.9pt] (\j,\j)--(\j+11,\j);}
\end{scope}\end{figC}
\begin{enumerate}[label=\alph*),leftmargin=1.6em,itemsep=1pt]
\item How many parallelograms are there in the figure altogether?
\item How many parallelograms contain the shaded small parallelogram?
\item How many parallelograms are made of \emph{exactly} $4$ small parallelograms?
\end{enumerate}
\end{probbox}
\probsol{
\textbf{Understanding the problem.} The figure has $4$ horizontal lines and $6$ slanted lines. As in Problem 27,
each parallelogram corresponds to a choice of $2$ horizontal lines (top and bottom sides) and $2$ slanted lines
(left and right sides).
\begin{recallbox}The number of ways to choose $2$ lines out of $k$ lines is $\dfrac{k(k-1)}{2}$.\end{recallbox}

\textbf{a)}
\step{1}{Count}
Choose $2$ of the $4$ horizontal lines: $\tfrac{4\cdot 3}{2}=6$ ways; choose $2$ of the $6$ slanted lines:
$\tfrac{6\cdot5}{2}=15$ ways. Number of parallelograms: $6\times 15=90$.

\textbf{b)}
\step{2}{Condition for containing the shaded cell}
Number the horizontal lines from the bottom up as $0,1,2,3$; the shaded cell lies between lines $1$ and $2$. A
parallelogram contains this cell when its bottom side is line $0$ or $1$ ($2$ ways) and its top side is line $2$ or
$3$ ($2$ ways). Number the slanted lines from left to right as $0,1,\dots,5$; the shaded cell lies between lines $2$
and $3$. The left side is line $0$, $1$ or $2$ ($3$ ways); the right side is line $3$, $4$ or $5$ ($3$ ways).
\step{3}{Multiply}
$2\times 2\times 3\times 3=36$ parallelograms.
\begin{proofbox}Why multiply? The four choices (bottom, top, left, right) are independent: every combination gives
exactly one parallelogram containing the shaded cell, and different combinations give different
parallelograms.\end{proofbox}

\textbf{c)}
\step{4}{``Sizes'' made of exactly $4$ cells}
A parallelogram $h$ rows high and $w$ columns wide consists of $h\times w$ cells. We need $h\times w=4$ with
$h\le 3$, $w\le 5$:
\begin{itemize}[leftmargin=1.4em,itemsep=1pt]
\item $h=1$, $w=4$: choose the row ($3$ ways) and $4$ adjacent columns out of $5$ ($2$ ways): $3\cdot 2=6$;
\item $h=2$, $w=2$: choose $2$ adjacent rows ($2$ ways) and $2$ adjacent columns ($4$ ways): $2\cdot 4=8$;
\item $h=4$, $w=1$: impossible, since there are only $3$ rows.
\end{itemize}
Total: $6+8=14$.
\begin{warnbox}Do not forget the $2\times 2$ type! And there are no parallelograms $4$ rows high.\end{warnbox}
\verify{Part a) counted by rows: $1$ row high: $3\cdot 15=45$; $2$ rows high: $2\cdot 15=30$; $3$ rows high: $15$
(each band has $5+4+3+2+1=15$ choices of columns). Total $90$. It matches.}
\ans{a) $90$. \ b) $36$. \ c) $14$.}
}

% ---------------------------------------------------------------- 49
\begin{probbox}{49}{\stars{2}}
\textit{(Extension of Problems 28 and 29.)} A $4\times 5\times 6$ rectangular box is built from $120$ small cubes with
edge $1$. All the outside faces of the box are painted, and then the small cubes are separated.
\begin{enumerate}[label=\alph*),leftmargin=1.6em,itemsep=1pt]
\item How many small cubes have exactly $3$ painted faces? exactly $2$? exactly $1$? no painted face at all?
\item What is the total number of painted small faces ($1\times 1$ squares)? Explain why the result equals the
surface area of the box.
\item An $n\times n\times n$ cube ($n\ge 3$) is painted and separated in the same way. Find $n$ such that the number
of small cubes with exactly $1$ painted face equals the number of small cubes with no painted face.
\end{enumerate}
\end{probbox}
\probsol{
\textbf{Understanding the problem.} The number of painted faces of a small cube depends on its \emph{position}: at a
corner, on an edge, on a face, or completely inside. Look at the front face (size $6\times 4$) to see this clearly:
\begin{figC}\begin{scope}[x=0.5cm,y=0.5cm]
\foreach \x in {0,...,5}{\foreach \z in {0,...,3}{
  \pgfmathtruncatemacro\ex{(\x==0||\x==5)?1:0}\pgfmathtruncatemacro\ez{(\z==0||\z==3)?1:0}
  \pgfmathtruncatemacro\t{\ex+\ez}
  \ifnum\t=2 \fill[tmOrange!60] (\x,\z) rectangle (\x+1,\z+1);\fi
  \ifnum\t=1 \fill[tmGreen!40] (\x,\z) rectangle (\x+1,\z+1);\fi
  \ifnum\t=0 \fill[tmBlue!30] (\x,\z) rectangle (\x+1,\z+1);\fi}}
\draw[Ink] (0,0) grid[step=1] (6,4);
\foreach \x in {0,...,6}{\draw[Ink] (\x,4)--(\x+2.5,5.75);}
\foreach \d in {0,...,5}{\draw[Ink] (0.5*\d,4+0.35*\d)--(6+0.5*\d,4+0.35*\d);}
\foreach \z in {0,...,4}{\draw[Ink] (6,\z)--(8.5,\z+1.75);}
\foreach \d in {0,...,5}{\draw[Ink] (6+0.5*\d,0.35*\d)--(6+0.5*\d,4+0.35*\d);}
\draw[Ink,line width=1.1pt] (0,0)--(6,0)--(8.5,1.75)--(8.5,5.75)--(2.5,5.75)--(0,4)--cycle (0,4)--(6,4)--(6,0) (6,4)--(8.5,5.75);
\node[below,font=\footnotesize] at (3,0){$6$};\node[left,font=\footnotesize] at (0,2){$4$};\node[right,font=\footnotesize] at (7.3,0.6){$5$};
\begin{scope}[xshift=5.2cm,yshift=0.4cm,font=\footnotesize]
\fill[tmOrange!60] (0,3) rectangle (0.6,3.6);\node[right] at (0.7,3.3){corner: $3$ painted faces};
\fill[tmGreen!40] (0,2) rectangle (0.6,2.6);\node[right] at (0.7,2.3){on an edge: $2$ faces};
\fill[tmBlue!30] (0,1) rectangle (0.6,1.6);\node[right] at (0.7,1.3){inside a face: $1$ face};
\end{scope}
\end{scope}\end{figC}

\textbf{a)}
\step{1}{Cubes at the corners: $3$ painted faces}
The box has $8$ corners, each with exactly one cube: $8$ cubes.
\step{2}{Cubes on an edge (not at its ends): $2$ painted faces}
The box has $12$ edges: $4$ of length $4$, $4$ of length $5$, $4$ of length $6$. An edge of length $L$ contains $L$
cubes; removing the $2$ cubes at its ends (already counted as corners) leaves $L-2$ cubes:
\[ 4\cdot(4-2)+4\cdot(5-2)+4\cdot(6-2)=8+12+16=36. \]
\step{3}{Cubes in the middle of a face: $1$ painted face}
On a $p\times q$ face, removing the outer border leaves $(p-2)(q-2)$ squares. The box has $2$ faces $4\times5$, $2$
faces $4\times 6$ and $2$ faces $5\times 6$:
\[ 2\cdot(2\cdot 3)+2\cdot(2\cdot 4)+2\cdot(3\cdot 4)=12+16+24=52. \]
\step{4}{Cubes completely inside: not painted}
Peel off a ``shell'' of thickness $1$ on every side; what remains is a box of
$(4-2)\times(5-2)\times(6-2)=2\times 3\times 4=24$ cubes.
\verify{$8+36+52+24=120$ -- exactly the total number of cubes.}

\textbf{b)}
\step{5}{Count the painted small faces}
$8\cdot 3+36\cdot 2+52\cdot 1+24\cdot 0=24+72+52=148$.
\begin{recallbox}The surface area of an $a\times b\times c$ box is $2(ab+bc+ca)$.\end{recallbox}
$2(4\cdot5+5\cdot6+6\cdot4)=2(20+30+24)=148$. The two results agree because each $1\times 1$ square on the outside
of the box is exactly one painted small face of exactly one cube -- we are simply counting the same collection in
two ways.

\textbf{c)}
\step{6}{Write the equation}
For an $n\times n\times n$ cube: each face has $(n-2)^2$ middle squares, so the $6$ faces give $6(n-2)^2$ cubes
with exactly $1$ painted face; the inside is a cube of $(n-2)^3$ small cubes. We need:
\[ (n-2)^3=6(n-2)^2. \]
\step{7}{Solve}
Since $n\ge 3$, $(n-2)^2>0$; dividing both sides by $(n-2)^2$: $n-2=6$, i.e.\ $n=8$.
\begin{warnbox}The condition $n\ge 3$ matters: for $n=2$ both numbers are $0$ (``equal'' in a trivial way), and we
must not divide by $0$.\end{warnbox}
\verify{$n=8$: $6\cdot 6^2=216$ and $6^3=216$. Correct.}
\ans{a) $8$; $36$; $52$; $24$ cubes. \ b) $148$. \ c) $n=8$.}
}

% ---------------------------------------------------------------- 50
\begin{probbox}{50}{\stars{2}}
\textit{(Extension of Problems 19 and 25.)} Ten unit squares (side $1$) are joined to form a flat shape such that:
adjacent squares share a whole side, and the shape is ``in one piece'' (from any square one can reach any other
square through shared sides). The perimeter of the shape is the total length of the square sides that are not
covered by another square.
\begin{figC}\begin{scope}[x=0.38cm,y=0.38cm]
\foreach \c in {(0,0),(1,0),(2,0),(2,1),(3,1),(4,1),(4,2),(1,1),(4,0),(5,1)}{\fill[tmBlue!30,draw=Ink] \c rectangle ++(1,1);}
\node[font=\footnotesize] at (3,-0.8){an example};
\end{scope}\end{figC}
\begin{enumerate}[label=\alph*),leftmargin=1.6em,itemsep=1pt]
\item Show that the perimeter is always an even number.
\item What is the greatest possible perimeter?
\item What is the smallest possible perimeter?
\item Which numbers can be the perimeter of such a shape?
\end{enumerate}
\end{probbox}
\probsol{
\textbf{Understanding the problem.} Let $k$ be the number of \emph{shared sides} (sides where two squares are joined).
Key idea: each shared side removes two square sides from the perimeter.

\textbf{a)}
\step{1}{A formula for the perimeter}
$10$ squares have $10\cdot 4=40$ sides in total. Each shared side is a side of two squares and is not on the
perimeter; every other square side is on the perimeter. Therefore
\[ P=40-2k. \]
$40$ and $2k$ are both even, so $P$ is even.

\textbf{b)}
\step{2}{A shape in one piece has at least $9$ shared sides}
\begin{proofbox}Rebuild the shape by placing one square at a time: start with any square, and each time add a
square next to one already placed (this is always possible because the shape is in one piece). Each added square
creates at least $1$ new shared side. There are $9$ additions, so $k\ge 9$.\end{proofbox}
\step{3}{Conclusion}
$P=40-2k\le 40-18=22$. Equality holds for a straight row $1\times 10$: perimeter $2(10+1)=22$. So the greatest
perimeter is $22$.

\textbf{c)}
\step{4}{Compare with the enclosing rectangle}
Let the smallest rectangle (with sides along the grid) containing the whole shape be $a\times b$ ($a$ columns,
$b$ rows).
\begin{proofbox}Every column of the enclosing rectangle contains at least one square. The highest square in the
column has its top side on the perimeter, and the lowest square has its bottom side on the perimeter: each column
contributes at least $2$ units of horizontal sides. Likewise each row contributes at least $2$ units of vertical
sides. So $P\ge 2a+2b=2(a+b)$.\end{proofbox}
\step{5}{Bound $a+b$}
The enclosing rectangle contains $10$ squares, so $ab\ge 10$. If $a+b\le 6$ then $ab\le 9$ (pairs with sum $6$:
$1\cdot5=5$, $2\cdot4=8$, $3\cdot3=9$; smaller sums give even smaller products) -- impossible. So $a+b\ge 7$ and
$P\ge 14$.
\step{6}{A shape with perimeter $14$ exists}
The $2\times 5$ rectangle: perimeter $2(2+5)=14$. So the smallest perimeter is $14$.

\textbf{d)}
\step{7}{The possible values}
From a), b), c): $P$ is even and $14\le P\le 22$, so $P\in\{14,16,18,20,22\}$. Each value is achieved by the
following ``two-row'' shapes (using the idea of Problem 25: the boundary is a staircase, so the perimeter equals that
of the enclosing rectangle):
\begin{figC}\begin{scope}[x=0.3cm,y=0.3cm]
\foreach \top/\bot/\xs/\lab in {5/5/0/14,6/4/6.5/16,7/3/14/18,8/2/22.5/20}{
  \foreach \i in {1,...,\top}{\fill[tmBlue!30,draw=Ink] (\xs+\i-1,1) rectangle ++(1,1);}
  \foreach \i in {1,...,\bot}{\fill[tmBlue!30,draw=Ink] (\xs+\i-1,0) rectangle ++(1,1);}
  \node[font=\footnotesize] at (\xs+\top/2,-0.9){$P=\lab$};}
\foreach \i in {1,...,10}{\fill[tmGreen!35,draw=Ink] (\i-1,-4) rectangle ++(1,1);}
\node[font=\footnotesize,anchor=west] at (10.5,-3.5){$P=22$};
\end{scope}\end{figC}
Top row $5,6,7,8$ squares, bottom row $5,4,3,2$ squares: the enclosing rectangles are $5\times2$, $6\times 2$,
$7\times 2$, $8\times 2$, with perimeters $14,16,18,20$; and the straight row gives $22$.
\begin{warnbox}``Compact'' does not automatically mean a smaller perimeter: a $3\times 3$ square with one extra square
sticking out has perimeter $12+2=14$, the same as the $2\times 5$ rectangle, not less.\end{warnbox}
\verify{The $6+4$ shape: $k=5+3+4=12$, $P=40-24=16$. The $8+2$ shape: $k=7+1+2=10$, $P=40-20=20$. Correct.}
\ans{a) $P=40-2k$ is even. \ b) $22$. \ c) $14$. \ d) $14$, $16$, $18$, $20$, $22$.}
}
